% Partie : théorie de la démonstration
% Chapitre : recherche de preuves
% Section : algorithme

\documentclass[../../../include/open-logic-section]{subfiles}

\begin{document}

\olfileid[fr]{pt}{ps}{sal}

\olsection{Algorithme de recherche pour \Log{G3c}}

Nous décrivons un algorithme de recherche pour \Log{G3c}.
Ce n'est ni le seul possible ni le plus efficace. Il est
cependant complet : si $\Gamma \Sequent \Delta$ possède
une preuve, il finit par la trouver et s'arrête.
S'il s'arrête sans succès, ou ne s'arrête jamais, alors
$\Gamma \Sequent \Delta$ est invalide. Ce dernier résultat
sera établi dans la section suivante.

L'algorithme doit garantir que toute formule non atomique
de $\Gamma \Sequent \Delta$, ou d'un séquent produit au
cours de la recherche, soit \emph{réduite} : elle doit
devenir la formule principale d'une règle appliquée
à rebours, autant de fois que nécessaire. Cette exigence
est celle de l'\emph{équité}.

La recherche procède par étapes. Chaque étape produit
un arbre de séquents. À l'étape suivante, on réduit une
formule dans chaque séquent supérieur dont la recherche
n'est pas déjà close.

À l'étape $0$, on part de $\Gamma \Sequent \Delta$.
Pour assurer l'équité, on attribue aux occurrences de
formules de $\Gamma$ et $\Delta$ des indices numériques
distincts, par exemple en les comptant de gauche à droite.
Ces indices sont écrits en exposant. Pour chercher une
preuve de $!A, !B \Sequent !C$, on peut ainsi partir,
à l'étape~$0$, de $!A^1, !B^2 \Sequent !C^3$.
Le plus grand exposant~$n$ d'un séquent indexé est
son \emph{indice maximal} : $3$ dans cet exemple.
Deux occurrences d'une même formule reçoivent elles
aussi des indices distincts.

En présence de quantificateurs, il faut également choisir
les termes employés pour réduire $\lexists[x][!A(x)]$
à droite ou $\lforall[x][!A(x)]$ à gauche.
Il suffit de termes construits à partir des constantes
$\Obj c_0$, $\Obj c_1$, $\Obj c_2$, \dots, et des symboles
de fonction présents dans~$\Gamma$ et $\Delta$.
On fixe une énumération effective de ces termes, dans
laquelle chaque terme a un rang fini, et une énumération
effective des constantes. On peut ainsi choisir la première
constante absente de $\Gamma \Sequent \Delta$.
À chaque séquent indexé est associé un ensemble
de constantes. À l'étape~$0$, il contient toutes
les constantes du séquent initial s'il y en a,
et vaut $\{\Obj c_0\}$ sinon.

Supposons qu'un arbre de séquents indexés, munis de
leurs ensembles de constantes, soit construit à
l'étape~$n$. À l'étape~$n+1$, on traite chacun des
séquents supérieurs comme suit.

Si une même formule~$!A$ figure à gauche et à droite,
ou si le membre gauche contient~$\lfalse$, on clôt
la recherche sur cette feuille et on ne la développe
plus. Un tel séquent possède une preuve dans \Log{G3c}.
La clôture sur une formule composée demandera de
greffer une preuve d'identité lors de la reconstruction
de la preuve finale.

Si la feuille n'est pas close et ne contient aucune
formule non atomique, l'algorithme s'arrête sans succès.
Le séquent initial $\Gamma \Sequent \Delta$ n'a pas
de preuve, comme l'établira le contre-modèle de la
section suivante.

Sinon, on choisit la formule non atomique~$!A$
d'indice~$i$ minimal. Le séquent considéré est alors
$!A^i, \Pi \Sequent \Lambda$ ou $\Pi \Sequent \Lambda, !A^i$.
On note $k$ le premier indice neuf, égal à l'indice
maximal du séquent augmenté de un, et $C$ l'ensemble
de constantes associé à ce séquent.

Réduire $!A^i$ consiste à engendrer les prémisses
de la règle déterminée par son opérateur principal
et par le membre, gauche ou droit, où elle figure.

\begin{enumerate}
  \item Si $!A \ident !B \land !C$ figure à gauche,
  on ajoute une seule prémisse au-dessus de
  $(!B \land !C)^i, \Pi \Sequent \Lambda$ :
  \[
  \Axiom$!B^k, !C^{k+1}, \Pi \fCenter \Lambda$
  \RightLabel{\LeftR{\land}}
  \UnaryInf$(!B \land !C)^i, \Pi \fCenter \Lambda$
  \DisplayProof
  \]
  L'ensemble de constantes de cette prémisse reste~$C$.

  \item Si $!A \ident !B \land !C$ figure à droite,
  on ajoute deux prémisses au-dessus de
  $\Pi \Sequent \Lambda, (!B \land !C)^i$ :
  \[
  \Axiom$\Pi \fCenter \Lambda, !B^k$
  \Axiom$\Pi \fCenter \Lambda, !C^k$
  \RightLabel{\RightR{\land}}
  \BinaryInf$\Pi \fCenter \Lambda, (!B \land !C)^i$
  \DisplayProof
  \]
  L'ensemble de constantes de chaque prémisse reste~$C$.

  \item Les cas $!A \ident \lnot !B$, $!A \ident !B \lor !C$
  et $!A \ident !B \lif !C$ sont analogues.
  Leur mise en forme détaillée est laissée en exercice ;
  les choix de prémisses sont précisés ci-dessous.

  \item Si $!A \ident \lexists[x][!B(x)]$ figure à gauche,
  on ajoute une prémisse au-dessus de
  $\lexists[x][!B(x)]^i, \Pi \Sequent \Lambda$ :
  \[
  \Axiom$!B(c)^k, \Pi \fCenter \Lambda$
  \RightLabel{\LeftR{\lexists}}
  \UnaryInf$\lexists[x][!B(x)]^i, \Pi \fCenter \Lambda$
  \DisplayProof
  \]
  où $c$ est la première constante absente de~$C$.
  L'ensemble de constantes de la prémisse est $C \cup \{c\}$.

  \item Si $!A \ident \lexists[x][!B(x)]$ figure à droite,
  on ajoute une prémisse au-dessus de
  $\Pi \Sequent \Lambda, \lexists[x][!B(x)]$ :
  \[
  \Axiom$\Pi \fCenter \Lambda, \lexists[x][!B(x)]^k, !B(t)^{k+1}$
  \RightLabel{\RightR{\lexists}}
  \UnaryInf$\Pi \fCenter \Lambda, \lexists[x][!B(x)]^i$
  \DisplayProof
  \]
  Le terme $t$ est le premier terme de l'énumération
  dont les constantes appartiennent à~$C$ et qui n'a
  pas encore été employé pour cette occurrence de
  $\lexists[x][!B(x)]$ à droite, sur la branche considérée.
  Si tous ces termes ont déjà été employés, on reprend
  la première constante de~$C$. Cet ensemble est
  toujours non vide. L'ensemble de constantes associé
  à la prémisse reste $C$.

  \item Les cas $!A \ident \lforall[x][!B(x)]$ sont
  analogues et leur mise en forme détaillée est laissée
  en exercice.
\end{enumerate}

Les autres cas se lisent dans les règles de \Log{G3c}.
Les contextes sont conservés. Sauf dans les deux cas
quantifiés réutilisables, on retire l'occurrence principale
et on ajoute les formules actives dans les prémisses.

Une négation $\lnot !B$ à gauche ajoute $!B$ à droite ;
à droite, elle ajoute $!B$ à gauche.
Une disjonction $!B \lor !C$ à gauche donne deux
prémisses, avec respectivement $!B$ et $!C$ à gauche ;
à droite, elle ajoute les deux formules à droite.
Une implication $!B \lif !C$ à gauche donne deux
prémisses, avec $!B$ à droite dans la première
et $!C$ à gauche dans la seconde ; à droite,
elle ajoute $!B$ à gauche et $!C$ à droite dans
une seule prémisse.
Une universelle $\lforall[x][!B(x)]$ à gauche est
conservée et réindexée, et l'on ajoute $!B(t)$
à gauche, en choisissant le terme et en conservant
son historique comme pour l'existentielle à droite.
À droite, on la remplace par $!B(c)$, avec la
première constante absente de $C$, que l'on ajoute
à l'ensemble associé.

Dans chaque prémisse, une seule nouvelle occurrence
reçoit l'indice $k$ ; lorsqu'il y en a deux, elles
reçoivent, dans l'ordre où elles viennent d'être
indiquées, $k$ et $k+1$. Pour les quantificateurs
conservés, l'occurrence réindexée reçoit $k$ et
l'instance reçoit $k+1$. Les autres indices sont
inchangés. Les ensembles de constantes restent
inchangés sauf lors de l'introduction d'une constante
propre. L'absurde est ici une formule atomique
du point de vue de la réduction : à gauche,
il clôt la feuille ; à droite, il n'est pas réduit.

Chaque occurrence quantifiée réutilisable possède,
sur sa branche, l'historique des termes déjà employés.
Cet historique est transmis à l'occurrence réindexée.
Lors d'une ramification, les occurrences du contexte
et leurs historiques sont copiés dans les deux
prémisses. Une nouvelle occurrence quantifiée
commence avec un historique vide. Sans symbole de fonction
d'arité positive, les termes disponibles sont les constantes
de~$C$ et l'on peut vérifier si elles ont toutes été utilisées.
Avec un tel symbole, on peut former une infinité de termes
sur~$C$ : l'historique étant fini, le parcours finit par
trouver un terme disponible encore inutilisé. Le rang fini
de chaque terme assure l'équité lorsque~$C$ croît.

La clôture sur une formule composée utilise la
\olref[pt][seq][ex]{prop:G3c-axioms} : ce séquent
d'identité est prouvable, même lorsqu'il n'est pas
lui-même un axiome atomique. On greffe cette preuve
à la feuille correspondante pour obtenir un arbre
%
de preuve de \Log{G3c}.

Si, à l'étape $n+1$, aucun séquent supérieur n'a
été développé, toutes les feuilles sont closes :
l'algorithme s'arrête avec succès. On obtient
une preuve de~$\Gamma \Sequent \Delta$ en effaçant
les indices de l'arbre de l'étape $n+1$ et en
greffant les preuves d'identité nécessaires.
Les feuilles de recherche ont toutes la forme
$!A, \Pi \Sequent \Lambda, !A$ ou
$\lfalse, \Pi \Sequent \Lambda$.
Après ces greffes, les feuilles sont des axiomes
et toutes les inférences sont des instances correctes
des règles de \Log{G3c}. C'est immédiat pour les
règles propositionnelles. À chaque étape, l'ensemble~$C$
associé à un séquent contient toutes les constantes
de ce séquent. Pour \LeftR{\lexists} et \RightR{\lforall},
la constante propre~$c$ est choisie hors de~$C$ :
elle est donc absente du séquent conclusion, ce
qui vérifie la condition de constante propre.
Nous avons ainsi le résultat suivant.

\begin{prop}
L'algorithme de recherche pour \Log{G3c} est correct :
s'il s'arrête avec succès, le séquent initial
$\Gamma \Sequent \Delta$ possède une preuve.
\end{prop}

\end{document}
